\subsection{An exact local exponent and infinitely many obstructed cubic fibers}

\begin{statement}{E.43}{Proved}{Every local class-two exponent and an infinite family excluded at higher stages}
For every prime $\mathfrak q\nmid3$ of
$K=\mathbb Q(\zeta)$, put $q=N\mathfrak q$ and
$s=v_3(q-1)$. The local rank-$44$ map of E.37 has
the exact cokernel exponent
\[
       \exp\operatorname{coker}d_{\mathfrak q}=3^{s+1}.
 \tag{E43.1}
\]
This is the exponent in the actual integral source
$M/L$ and target $M_0$ of (E37.6).

Let $\mathcal P$ be the following set of rational primes:
\[
 \mathcal P=\{p:p\equiv1\pmod9,\quad
                         2^{(p-1)/3}\not\equiv1\pmod p\}.
 \tag{E43.2}
\]
It has natural density $1/9$, in particular it is
infinite. Use the actual rational tangential class
$\alpha_P$ and rational peripheral vector $e_Q$
of E.38. For $p\in\mathcal P$ and every integer
$k\ge v_3(p-1)+1$, the class
\[
 \alpha_{p,k}=\alpha_P+3^k\kappa_3(p)e_Q,
       \qquad\text{toric coordinates }(2,p^{3^k}),
 \tag{E43.3}
\]
has a nonempty full intrinsic class-two fiber, with
the complete parameter group (E37.12). Every section
in that fiber fails to lift to $N/N''$, to $N$,
or to any marked quotient in (E38.2). In fact the
local $H$ obstruction excludes those lifts for
every $k\ge0$, independently of the positive
class-two bound. The classes for distinct primes
are distinct already in the minimum torus quotient.

The primes in (E43.2) through $1000$ are
\[
\begin{gathered}
19,37,73,163,181,199,271,379,487,523,541,577,\\
613,631,757,829,883,937,991.
\end{gathered}
\tag{E43.4}
\]
All $38$ local rank-$44$ matrices at their two
primes in $K$ have been computed and agree with
(E43.1). For $p=19$ the bound is $k=3$, recovering
the finite class-three obstruction of E.38.

\emph{Manuscript: \S51.16, folios 733v--736v, scans 172--178;
Part E(c),(d) and the fourth exactness follow-up.
The positive calculation uses the full intrinsic
matrix, and the negative calculation uses the
actual quotient maps (E38.9).}
\end{statement}
\begin{proof}
We first prove the exact exponent. Write
$\mathcal O=\mathbb Z_3[\zeta]$ and
$z=\psi_{\mathfrak q}-1$. The scalar is primary,
so $z\in3\mathcal O$, and $z\ne0$ since its
Frobenius norm is $q>1$. The two Smith exponents of
multiplication by $z$ on $\mathcal O$ are
\[
 m=\lfloor v_{1-\zeta}(z)/2\rfloor,\qquad
 h=\lceil v_{1-\zeta}(z)/2\rceil.
 \tag{E43.5}
\]
Indeed $\mathcal O$ is a discrete valuation ring,
$3$ is a unit times $(1-\zeta)^2$, multiplication
by a unit is an integral automorphism, and
multiplication by $1-\zeta$ has Smith factors
$1,3$. These facts give (E43.5) by separating
the even and odd powers of the uniformizer.

The trace of $\mathcal O$ is $\mathbb Z_3$ and
the trace of $(1-\zeta)\mathcal O$ is $3\mathbb Z_3$:
use the bases $1,\zeta$ and $1-\zeta,\zeta-\zeta^2$.
Consequently $\operatorname{Tr}(z)\in3^h\mathbb Z_3$,
and $N(z)\in3^{m+h}\mathbb Z_3$. The identity
\[
 q-1=N(1+z)-1=\operatorname{Tr}(z)+N(z)
\]
therefore gives $s\ge h$.

Inside $M_0$ consider the explicit stable sublattice
\[
 S_0=\bigwedge^2T_3T_D+
                  (T_3T_D\wedge j(T_3E)).
 \tag{E43.6}
\]
For the elliptic basis $t_i$, the columns
$j(t_i)=3f_i+U_i$ in E.37 show that
$3M_0\subset S_0$: the toric wedges are already
in $S_0$, and three times each mixed wedge is a
mixed $j$-wedge minus toric wedges. The inclusion
$S_0\to M/L$ is injective, since $L$ has nonzero
$f_1\wedge f_2$ coordinate while $S_0$ has none.
On its toric wedges, $d_{\mathfrak q}$ is
multiplication by $q-1$. On its mixed wedges it
is the actual matrix $A_E-1$, whose largest
Smith exponent is $h\le s$. Hence
$3^s d_{\mathfrak q}^{-1}S_0\subset S_0$.
Together with $3M_0\subset S_0$ this proves
\[
             3^{s+1}M_0\subset d_{\mathfrak q}(M/L).
 \tag{E43.7}
\]

We establish equality of the exponent by an integral
direct summand. The two columns of $U$ span a primitive
rank-two sublattice of $\mathbb Z_3^8$: their minor
at $A_0,A_2$ is a unit. A change of torus basis sends
$U$ to $\binom{I_2}{0}$. In this basis the actual
Frobenius module decomposes as
\[
 V=V_4\oplus T_6,\qquad
 A_{V_4}=\begin{pmatrix}A_E&0\\C&qI_2\end{pmatrix},
 \quad C=(A_E-qI_2)/3,
 \tag{E43.8}
\]
where $T_6$ is toric of rank six and $L\subset
\bigwedge^2V_4$. Thus $V_4\otimes T_6$ is an
actual direct summand of both $M/L$ and $M_0$,
and $d_{\mathfrak q}$ acts there as six copies
of $A_{V_4}-1$.

Identify its two-dimensional blocks with
$\mathcal O$. The element represented by $C$ is
$-\psi_{\mathfrak q}\overline z/3$.
Since conjugation preserves the uniformizer
valuation, the matrix
\[
       (A_E-1)C^{-1}
            =-3\psi_{\mathfrak q}^{-1}z/\overline z
\]
is three times an integral automorphism. Also
$(q-1)C^{-1}$ is integral: its uniformizer
valuation is $2s+2-v_{1-\zeta}(z)\ge2$.
Integral block row and column operations on
\[
 A_{V_4}-1=\begin{pmatrix}A_E-1&0\\C&(q-1)I_2\end{pmatrix}
\]
therefore give a block diagonal matrix with
blocks $C$ and three times $q-1$ times an
integral automorphism. Its Smith exponents are
\[
                 m-1,\ h-1,\ s+1,\ s+1.
 \tag{E43.9}
\]
This direct summand has exponent $3^{s+1}$.
Combined with (E43.7) it proves (E43.1) for
every prime under consideration.

We next verify the positive global lifting assertion.
The class $\alpha_P$ comes from a section of the
full arithmetic group and is unramified outside
$2,3$. At the prime of $K$ above two, an odd
rational prime $p$ lies in $1+2\mathcal O_{K,2}$.
Cubing is an automorphism of this pro-two unit
group, so $\kappa_3(p)$ vanishes there. The
local class in (E43.3) thus has the same lifting
property as $\alpha_P$ at two.

At either prime above $p$, normalize the tame
valuation of $p$ to one. With the notation of
E.37 the full local numerator is
\[
 D_{\mathfrak p}
       =3^k(a_P-k_{\mathfrak p,Q})\wedge e_Q
                   \in3^k M_0.
 \tag{E43.10}
\]
The self-wedge term from the unramified value of
$\kappa_3(p)$ is zero. Since $k\ge s+1$, (E43.1)
makes the unique correction integral at both
primes. All other primes away from three have
zero inertia value, apart from the already
liftable prime above two. The correction family
has finite support and satisfies the full
continuity condition (E37.11). E.37 proves the
claimed global class-two lift and identifies
every class in its fiber.

For the negative assertion, use the unit tangent
at $P$ and the quotient $N\to H$ of E.38.
At $p\in\mathcal P$ the first inertia coordinate
is $a=0$, and the inertia augmentation is
$c=3^{k+1}\ne0$. The first Frobenius coordinate
is $b=3\kappa_3(2)(F)$, up to addition of
$3(p-1)\mathbb Z_3$ by geometric conjugation.
The residue condition in (E43.2) says precisely
that $\kappa_3(2)(F)$ is a unit modulo three.
Thus $v_3(b)=1$, while $v_3(p-1)\ge2$.
Condition (E42.3) fails for every representative
and every nonzero $c$. This excludes a local
$H$-section above the prescribed abelian class.
Any global section in any of the stated higher
quotients would restrict and map to that
excluded local section. The exact maps are
(E38.9), and $N$ maps onto each of its quotients.
This proves the nonlifting statements through
those morphisms.

Finally let $L_9=\mathbb Q(\zeta_9,2^{1/9})$.
At a prime above two, the cyclotomic field
$\mathbb Q(\zeta_9)$ is unramified and
$X^9-2$ is Eisenstein. The Kummer degree is
therefore nine, and $[L_9:\mathbb Q]=54$.
An automorphism is the actual pair
\[
 (u,v):\quad\zeta_9\mapsto\zeta_9^u,
       \quad2^{1/9}\mapsto\zeta_9^v2^{1/9},
 \qquad (u,v)(u',v')=(uu',v+uv').
\]
The conjugacy class of $(1,1)$ consists of
the six pairs $(1,v)$ with $v\in(\mathbb Z/9)^\times$.
At an unramified rational prime, this class
is exactly the two residue requirements in
(E43.2). Chebotarev's theorem
\cite[V, Theorem 3.23]{E_MilneCFT}, applied to
this explicit finite Galois extension, gives
density $6/54=1/9$.
The toric valuation at $p$ in (E43.3) is
$3^k$, whereas the same coordinate for a
different prime has valuation zero at $p$.
This proves distinctness through the actual
minimum torus quotient.

The program \texttt{e\_infinite\_obstructions.py}
enumerates every eligible prime through $1000$,
counts $E(\mathbb F_p)$ on the model
$y^2+y=x^3-7$, finds both primary CM scalars,
and computes both complete rank-$44$ Smith
lists. The exact matrices and lists are saved.
It checks (E43.1) in all $38$ cases and records
the first forced nonintegral coefficient from
(E42.10). The uniform exponent, infinitude,
and nonlifting at every $k$ are supplied by
the preceding proofs.
\end{proof}

